\section{Discussion}
The first question for a replay implementation is which state its displacement updates. A measured-origin controller exposes that state through the hand observation. For an accumulating controller the initial target, update law and timed sequence of applied increments suffice. A sampled target anchors that reconstruction at its sample cycle. If later inputs are ambiguous, the uncertainty grows again before the packet arrives. An execution bit settles the ambiguity only to the extent that its producer identifies the applied input under the model in use. The local Franka experiment illustrates the value of an application-cycle field when its update law and parameters are known. DROID illustrates the difficulty of identifying an equivalent field in an archived physical demonstration.

A program can still decide to change a command with an uncertain origin. The guarantee in~\eqref{eq:guard} compares two displacements at the same origin and desired target for one delivered write. Even a certified improvement on that write may increase trajectory error, consume more reads or leave object placement unchanged. Synchronized physical replay with independently logged applied inputs and object outcomes would test how much these distinctions matter under real sensing and communication delays~\cite{li2025simpler}.

The experiments use scripted time-aligned sources and Cartesian translations. The origin set depends on the assumed fault model. A forwarding flag interpreted as proof of full application can exclude the true origin when inputs are attenuated. The quantized-hand experiments model deterministic coordinate rounding. The all-success confirmation cohort limits inference about uncommon rescues. Costs of reading also depend on whether the receiver already streams the target. These conditions delimit when a replay implementation can use the guarantees reported here.
